% !TeX root = ../papers/nc-p3.tex
\section{節点・尖点核の四次鉛筆と二項式の行列表示}
\label{sec:singular-elliptic-cofactor-pencil}
\begin{theorem}\label{thm:singular-elliptic-cofactor-pencil}
四辺有限な標準有限Frobenius窓の整な八余因子の場合、核四次曲線 $E$ の唯一の特異点を
$p$ とすると、非零二次式 $p^{\mathsf t}N(x)p$ は二つの一次式の積である。
節点では乗法群のパラメータ $t\sim t^{-1}$、尖点では加法群の $t\sim-t$ の
四次鉛筆があり、正規化の境界を加えた商の一点が一意に選ばれる。
節点の非自明二捩れ $t=-1$ と尖点の境界では上の二次式は平方である。
有限の非分岐元では、$N$ は以下の明示Bezout行列に限られる。
\end{theorem}
\begin{proof}[特異点での行列小行列式]
正規化の変数を $z$ とし、節点の二枝を $0,\infty$、尖点の枝を0とする。
$L=\mathcal O_E(1)$ の四切断は節点で
$(1+z^4,z,z^2,z^3)$、尖点で $(1,z^2,z^3,z^4)$ を基底にできる。
節点の非零な二端貼合係数は変数の倍率で吸収する。
尖点の条件は初め $a_1=k a_0$ であるが、0を固定する射影変換で二項四次式を移せば
一次係数に定数係数の任意倍を加えられる。$L^2$ についても同じ貼合条件を二乗する。
八切断に対する汎関数を $h$ とし、$h_i$ は該当する単項式 $z^i$ の値、
$h_*$ は節点で $1+z^8$ の値、尖点で1の値とする。
元の $L,L$ 対合に六次式座標を代入して $\operatorname{adj}N$ が得られるように規格化する。

節点で $P_t,Q_t$ は第一基底を $1+t z^4,1+t^{-1}z^4$ に取り替える。
その積の四行対合は
\[
 C_t=\begin{pmatrix}h_*+(t+t^{-1})h_4&(u+t v)^{\mathsf t}\\
 u+t^{-1}v&H\end{pmatrix},
\quad H=(h_{i+j})_{1\le i,j\le3},\quad
 u=(h_1,h_2,h_3)^{\mathsf t},\ v=(h_5,h_6,h_7)^{\mathsf t}.
\]
$\det C_t$ は $t+t^{-1}$ に関して一次で、最高係数は
$-\det(h_{1+i+j})_{0\le i,j\le3}$ である。四行Hankel行列の列を循環して
左上三行を $H$ に取れば分かる。尖点では第一基底を $1+tz,1-tz$ として
\[
 C_t=\begin{pmatrix}h_*-t^2 h_2&(u+t v)^{\mathsf t}\\u-t v&H\end{pmatrix},
\quad H=(h_{i+j})_{2\le i,j\le4},\quad
 u=(h_2,h_3,h_4)^{\mathsf t},\ v=(h_3,h_4,h_5)^{\mathsf t}.
\]
$t^2$ の最高係数は $-\det(h_{2+i+j})_{0\le i,j\le3}$ である。
いずれも最高係数は非零な多項式で $h_*$ を含まず、元の行列式とは独立である。

$D$ 上の一般点ではどの有限パラメータの対合も一点評価で階数1になる。
境界Hankel行列でも評価が共通因子を除き $(1,z,z^2,z^3)$ の平方になる。
よって代入した全ての四次行列式が $d^3$ の定数倍で、元のものは非零、
選ばれる鉛筆の零元は一意である。また $\det H=d^2 f$、
$f=p^{\mathsf t}N(x)p\ne0$ は余因子の相補小行列式公式である。
どの有限元でも $H$ は左上余因子であり、境界Hankel行列でも $H$ は三次小行列式である。
従って選ばれた元の一般階数はちょうど3である。

その随伴は $d^2$ で割れ、商は一般階数1の二次式行列である。
有限元には固定左右核がなく（非零切断と相手の切断の積は零でない）、
境界でも七つの独立座標を持つHankel行列に固定核はない。
原始分解で商は一次列の外積となる。従ってその成分である $f$ が二一次式の積である。
対称な有限分岐元では $H$ が主小行列なので平方となり、尖点境界でも $H$ は主小行列である。
節点境界では両端の別々の行と列を落とすため積である。
\end{proof}

\begin{proposition}\label{prop:singular-elliptic-binary-bezout}
有限非分岐の場合、入力を $a\in P_t=H^0(L\eta)$ の係数に取り、
$b=Ta\in Q_t=H^0(L\eta^{-1})$ とするある線型同型 $T$ がある。
このとき $N$ は $z,w$ の $L$ 基底で次の対称二変数式を読んだ行列である。
\[
\begin{array}{ll}
\text{節点}&
\displaystyle\frac{(z-tw)a(z)b(w)+(tz-w)a(w)b(z)}{z-w},\qquad t\ne 0,1,-1,\\[6pt]
\text{尖点}&
\displaystyle\frac{a(z)b(w)+a(w)b(z)}{t}
   +\frac{zw\{a(z)b(w)-a(w)b(z)\}}{z-w},\qquad t\ne0 .
\end{array}
\]
全体の非零定数は基底に吸収する。
\end{proposition}
\begin{proof}
先の一次列 $a(x),b(x)$ が $C_t$ の左右核の切断である。
それらの第一係数の積が $f\ne0$ なので一般には特異点で消えない。
一般の二切断に共通因子 $A$ があれば滑らかな部分の有効因子で、
$P_t\ne Q_t$ より次数は3以下となる。
$aH^0(Q_t)\cap bH^0(P_t)=abH^0(\mathcal O_E(A))$ であり、
正次数の線束の第一コホモロジー消滅から余因子汎関数が長さ3以下の $A$ での
評価を経由してしまう。元の可逆対合に反する。
よって一般には共通零点がなく、二積空間の和の零化直線が余因子を決定する。

表示した式は常に二つの変数について $L$ の貼合条件を満たす。
節点では $K=(z-tw)/(z-w)$ が二端で互いに逆の貼合係数を与え、
尖点では $K=1/t+zw/(z-w)$ の0での対数微分が二変数で $-t,t$ になる。
各式は $K a(z)b(w)$ とその転置の和で、対角の極が消える。
一般の $a$ の相異なる滑らかな四根で対角となり、逆は四評価の重み付き和である。
$aH^0(Q_t)$ に加えて $bH^0(P_t)$ を消すことは、前節と同じ留数公式で確認できる。
不変微分は節点で $dz/z$、尖点で $dz/z^2$ と取る。$h/a$ が $E$ 上正則な特異点では
二端の留数が相殺され、または尖点の一次係数が零である。
表示の行列式は正規化上の二四次式の結果式に非零定数倍で等しい。
実際滑らかな共通根を印した入射の像は既約で、片側の一般切断の四根が
相手に相異なる四超平面を課す。両者とも双次数 $(4,4)$ で、この既約因子を含む。
従って一般の可逆な対合の逆に一致した射影公式は互いに素な対で成立する。
両辺の二次性ともとの十成分の原始性で恒等的に定数比例する。
一次写像が核を持てば共通基点が生じるので、両方とも同型である。
\end{proof}

\path{computations/verify_singular_elliptic_cofactor_pencil.py} は二つの特異型で
核の貼合、二項結果式との倍率、境界のHankel行列式を検算する。有限非分岐パラメータと
同型 $T$ の各有理標本で二次成分は十独立だが、四次小行列式は十九独立である。
有限辺の二十点を通る四次空間は十五次元なので、この開条件の範囲は排除される。
